\section{Peramalan}
Peramalan deret waktu \textit{(time series foecasting)} merupakan komponen kunci yang digunakan dalam berbagai bidang untuk memprediksi kejadian di masa depan dan mengidentifikasi tren. 
Metode peramalan terbagi menjadi dua yaitu kualitataif (berbasis opini ahli) dan kuantitatif (berbasis analisis statistik dan matematis), dengan deret waktu sebagai salah satu pendekatan kuantitatif yang umum digunakan karena menganalisis data berurutan berdasarkan waktu \citep{petropoulos2021forecasting}. 
Model-model peramalan telah berkembang secara bertahap seiring waktu. 
Pada awalnya, metode prediksi deret waktu mengandalkan prosedur statistik murni seperti analisis regresi dan metode ARIMA \citep{putra2019sugar}. 
Metode statistik ini bekerja dengan baik untuk data dengan hubungan linier. 
Untuk data non linier, model kecerdasan buatan terutama jaringan saraf telah digunakan \citep{marino2016building}. 
Baru baru ini, metode hibrida yang menggabungkan metode statistik dengan pendeketan deep learning telah dikembangkan untuk menghasilkan prediksi yang lebih akurat \citep{jagtap2019rainfall}. 
Bahkan model konvensional seperti ARIMA sering kali gagal memebrikan prediksi curah hujan yang akurat ketika diterapkan pada wilayah dengan variasi geografis yang kompleks sehingga memerlukan model alternatif. 
Maka dari itu, dalam konteks penelitian ini menggunakan model yang efektif ketika berhadapan dengan data spatio temporal. 
Untuk mengetahui model itu efektif atau tidaknya dalam penelitian ini menggunakan evaluasi RMSE. 
Berikut adalah rumus umum untuk menghitung RMSE.

\begin{equation}\label{rmse}
    \text{RMSE} = \sqrt{ \frac{1}{n} \sum_{t=1}^{n} (y_t - \hat{y}_t)^2 }
\end{equation}
Persamaan~(\ref{rmse}) menunjukkan rumus \textit{Root Mean Square Error} (RMSE).
Keterangan:
\begin{itemize}
  \item $y_t$ : nilai aktual pada waktu ke-$t$
  \item $\hat{y}_t$ : nilai prediksi pada waktu ke-$t$
  \item $n$ : jumlah total data atau observasi yang diprediksi
\end{itemize}
Nilai-nilai seperti RMSE merupakan ukuran tingkat kesalahan prediksi. Semakin rendah nilai-nilai ini, maka semakin akurat model dalam memprediksi data.